\documentclass[10pt,a4paper]{amsart}
\usepackage[margin=19mm]{geometry}
\usepackage{amsmath,amssymb,mathtools}
\usepackage[scr=rsfs,cal=euler]{mathalfa}
\usepackage{xfrac}
\usepackage[unicode,hidelinks,bookmarks=false]{hyperref}
\newcommand{\ie}{i.e.\ }
\newcommand{\cf}{cf.\ }
\newcommand{\resp}{respectively\ }
\newcommand{\Subsection}{\subsection}
\providecommand{\nobreakdash}{\nobreak\hbox{-}}
\setlength{\emergencystretch}{2em}
\multlinegap0pt
\usepackage{amssymb}
\usepackage{tikz}
\usepackage[shortlabels]{enumitem}
\newtheorem{lemma}{Lemma}
\newtheorem{corollary}{Corollary}
\newcommand{\Z}{\mathbb{Z}}
\title{Borel Edge Colorings for Finite Dimensional Groups}
\author{Felix Weilacher}
\date{Source: arXiv:2104.14646v2 (CC BY 4.0)}
\begin{document}
\maketitle

\noindent\emph{Source and licence.} Borel Edge Colorings for Finite Dimensional Groups. Version arXiv:2104.14646v2 (CC BY 4.0), distributed by arXiv under the Creative Commons Attribution 4.0 International licence, http://creativecommons.org/licenses/by/4.0/, as recorded in the arXiv licence metadata for this exact version and shown on the abstract page https://arxiv.org/abs/2104.14646v2. Full licence notice: \texttt{licenses/arxiv-2104.14646-CC-BY-4.0.txt}.\par\smallskip

\noindent\textit{Editorial scope.} Counted unit: the article's lemma lem:quotient (source0.tex lines 347--349). The article's complete original proof of it is delivered with it (source0.tex lines 351--364, one \texttt{proof} environment of 14 lines). No mathematics is written, repaired or re-derived: the statement and the proof are the article's own lines, in the article's own notation, and no retained line is substituted or re-flowed.  Retained uncounted as the source-local context the proof needs: lines 188--190; lines 212--214.  Nothing was omitted from any retained span, so the package carries no omission record.  No retained line contains a citation, so the wrapper carries no bibliography and no bibliography entry.  No retained line contains a figure environment, an image-inclusion command or drawing code, so no image is delivered and the package declares no asset dependency.  Editorial formatting only, in addition: an AMS \texttt{amsart} preamble that restates the article's own declarations and macros used by the retained lines, namely the environments \texttt{lemma}, \texttt{corollary} on separate counters, together with the standard packages those lines need.  The retained environments print in this document as Lemma 1, Corollary 1 in order of appearance, whereas the article prints the same items with the numbering of its own full document.  The complete native source of the article as distributed in the arXiv e-print for this exact version is bundled unchanged as \texttt{sources/110892/source0.tex}.
\begin{lemma}\label{lem:Z}
Let $a:\Z \curvearrowright X$ be a free Borel action of $\Z$, and $S$ the usual generating set for $\Z$. Suppose $V \subset X$ is a Borel set with the property that every injective $G(a,S)$-ray contains infinitely many edges within $V$. Then there is a Borel 3-edge coloring of $G(a,S)$, say with the colors $1,2$, and $3$, such that the color 3 only occurs on edges contained within $V$.
\end{lemma}

\begin{corollary}\label{cor:d+1}
Let $a:\Gamma \curvearrowright X$ be a free Borel action of a marked group $(\Gamma,S)$ on a standard Borel space $X$ with Borel asymptotic separation index 1. Suppose $\Delta \leq \Gamma$ is a finite normal subgroup of $\Gamma$, and that for each $\gamma \in S \setminus \Delta$, the image of $\gamma$ in the quotient $\Gamma/\Delta$ does not have odd order. Then $\chi_B'(G(a,S)) \leq |S|+1$.
\end{corollary}

\begin{lemma}\label{lem:quotient}
Let $\Delta' \leq \Gamma$ be a finite proper subgroup. Let $\overline{\Gamma} = \Gamma/\Delta'$, and then define $\overline{a}:\overline{\Gamma} \curvearrowright \overline{X}$, $S_0$, and $\overline{S_0}$ as in the proof of Corollary \ref{cor:d+1}. If $\chi'(G(\overline{a},\overline{S_0})) = |{S_0}|$, then $\chi'(G(a,S)) = |S|$. Likewise for the Borel edge chromatic numbers of these graphs.
\end{lemma}

\begin{proof}
Let $S_1 = S \cap \Delta' = S \setminus S_0$. By hypothesis, we may find a $|S_0|$-edge coloring of $G(\overline{a},\overline{S_0})$, say using the set of colors $C$. As in the proof of Corollary \ref{cor:d+1}, we may lift this to a $|S_0|$-edge coloring of $G(a,S_0)$ by giving the edge $(x,\gamma \cdot x)$ for $x \in X$ and $\gamma \in S_0$ whichever color was given to the edge $(\Delta' \cdot x, \gamma \Delta' \cdot x)$ in $G(\overline{a},\overline{S_0})$ corresponding to $\gamma$.


Fix a color $c \in C$, which exists since $\Delta'$ is proper, so $S_0$ is nonempty. Consider the subgraph $G'$ of $G(a,S)$ consisting of $G(a,S_1)$ along with all the edges colored $c$ in the previous step. 

Since, in our coloring of $G(\overline{a},\overline{S_0})$, the set of edges colored $c$ gives a perfect matching, each connected component of our subgraph looks like the following: Two adjacent $\Delta$-orbits, say $y$ and $\gamma \cdot y$ for some $\gamma \in S_0$, with all the internal edges from $S_1$, along with all edges between them of the form $(x,\gamma \cdot x)$ for $x \in y$.

Fix a set $C'$ of $|S_1|+1$ colors disjoint from $C$. We now use $C'$ to edge color the above component. First, by Vizing's theorem, we may use $C'$ to color the edges within $y$. Now, for every such edge, say $(x,x')$, give the edge $(\gamma \cdot x, \gamma \cdot x')$ the same color. Since $\Gamma$ is abelian, this colors all edges within $\gamma \cdot y$ without conflict. Finally, for each $x \in y$, by this construction the set of colors of edges in $y$ meeting $x$ is the same as the set of colors of edges in $\gamma \cdot y$ meeting $\gamma \cdot x$. This set has size $|S_1|$, though, so we still have one free color, and can assign it to the edge $(x,\gamma \cdot x)$, completing our coloring.

Thus we can edge color $G'$ using $C'$, and we already have an edge coloring of $G(a,S) \setminus G'$ using $C \setminus \{c\}$, so unioning these gives a coloring using $|C|'+|C|-1 = |S_1|+1 + |S_0| - 1 = |S|$ colors as desired.

Finally, if our original coloring was Borel, our lift will still be Borel, and then our construction can be done in a Borel fashion since we only need to work with finite components. Again this uses the argument from Lemma \ref{lem:Z}.
\end{proof}

\end{document}
